\documentclass[11pt]{article}
\usepackage[utf8]{inputenc}
\usepackage[T1]{fontenc}
\usepackage[margin=1in]{geometry}
\usepackage{amsmath, amssymb}
\usepackage{booktabs}
\usepackage{array}
\usepackage{hyperref}

\begin{document}

% =============================
% Exam definition (included)
% =============================
\begin{center}
{\LARGE First Exam — Lessons 0 \& 1 (Version B)}\\[0.75em]
\end{center}

\noindent\textbf{Time allowed:} 60 minutes\\
\textbf{Resources:} Printed course notes and personal handwritten notes only. Electronic devices, calculators, and statistical tables are not permitted.\\
\textbf{Instructions:} Answer all four questions. Show algebraic steps and clearly label each part. Unless stated otherwise, provide exact values (fractions, radicals, exponentials). Partial credit is available when reasoning is clearly explained.

\vspace{0.75em}
\noindent\begin{tabular}{@{}p{2.2cm}p{8.2cm}r@{}}
\toprule
\textbf{Question} & \textbf{Topic} & \textbf{Points}\\
\midrule
1 & Probability structures and independence & 25\\
2 & Bayes' rule and conditional reasoning & 25\\
3 & Uniform model derivations & 25\\
4 & Interpreting summaries and convergence & 25\\
\midrule
\textbf{Total} & & \textbf{100}\\
\bottomrule
\end{tabular}

\vspace{1em}
\hrule
\vspace{1em}

\section*{Question 1 — Probability Foundations (25 points)}

A die is rolled twice. Throughout parts (a)–(c), assume the die is fair.

\begin{enumerate}
  \item (6 pts) Specify a probability space $\,(\Omega, \mathcal{F}, P)\,$ for this experiment. Clearly describe the sample space, a natural $\sigma$-algebra, and the probability measure.
  \item (6 pts) Let $A$ be the event ``the sum of the two rolls is at least 10'' and $B$ the event ``the first roll shows a 5 or 6.'' Compute $P(A)$, $P(B)$, and $P(A \cap B)$.
  \item (5 pts) Are $A$ and $B$ independent? Justify your answer using the formal definition.
  \item (8 pts) Now suppose the die is biased such that the probability of rolling a 6 is $q$ (where $0 < q < \tfrac{1}{6}$), while outcomes 1 through 5 are equally likely with total probability $1-q$. With the same definition of event $B$, derive an expression for $P_q(B)$ in terms of $q$. For which value of $q$ does $P_q(B) = \tfrac{1}{3}$? Provide algebraic reasoning.
\end{enumerate}

\vspace{0.75em}
\hrule
\vspace{0.75em}

\section*{Question 2 — Bayes' Rule in Fraud Detection (25 points)}

A bank uses an automated system to flag potentially fraudulent credit card transactions. Historical data indicate that $\tfrac{1}{50}$ of transactions are actually fraudulent. The system has sensitivity (true positive rate) $\tfrac{4}{5}$ and specificity (true negative rate) $\tfrac{39}{40}$.

\medskip
\noindent\textit{Recall:} sensitivity $= P(\text{Flag} \mid \text{Fraud})$ and specificity $= P(\text{No Flag} \mid \text{No Fraud})$. The false positive rate is $1-\text{specificity} = P(\text{Flag} \mid \text{No Fraud})$.

\begin{enumerate}
  \item (7 pts) Compute the probability that a randomly selected transaction is flagged as suspicious.
  \item (8 pts) Given that a transaction is flagged, determine the probability that it is actually fraudulent. Express your answer as a reduced fraction.
  \item (6 pts) Given that a transaction is not flagged, determine the probability that it is actually fraudulent. Express your answer as a reduced fraction.
  \item (4 pts) Let $\pi = P(\text{Fraud})$, $s = P(\text{Flag} \mid \text{Fraud})$ (sensitivity), and $c = P(\text{No Flag} \mid \text{No Fraud})$ (specificity). Derive a formula for $P(\text{Fraud} \mid \text{Flag})$ in terms of $\pi$, $s$, and $c$, and briefly state how increasing $c$ affects this probability.
\end{enumerate}

\vspace{0.75em}
\hrule
\vspace{0.75em}

\section*{Question 3 — Uniform Waiting-Time Model (25 points)}

Let $T$ be the time (in minutes) until the next customer arrives at a service counter, and suppose $T$ follows a uniform distribution on $[0,8]$.

\medskip
\noindent\textbf{Recall:} The probability density function of a $\mathrm{Uniform}(a,b)$ random variable is
\[
f_T(t) =
\begin{cases}
\dfrac{1}{b-a}, & a \le t \le b,\\
0, & \text{otherwise},
\end{cases}
\]
and the cumulative distribution function is $F_T(t) = 0$ for $t<a$, $F_T(t) = \dfrac{t-a}{b-a}$ for $a \le t \le b$, and $F_T(t) = 1$ for $t>b$.

\begin{enumerate}
  \item (5 pts) Specialise the CDF $F_T(t)$ for $T \sim \mathrm{Uniform}(0,8)$.
  \item (5 pts) Using your result from part (1), compute $P(2 \le T \le 6)$.
  \item (9 pts) Compute $E[T]$ and $\operatorname{Var}(T)$ directly from integrals. Show the essential integration steps leading to exact values.
  \item (6 pts) Define $Y = 2T + 5$. Compute $E[Y]$ and $\operatorname{Var}(Y)$ using properties of expectation and variance. Provide exact values.
\end{enumerate}

\vspace{0.75em}
\hrule
\vspace{0.75em}

\section*{Question 4 — Interpreting Summaries and Convergence (25 points)}

Table~\ref{tab:b-sim-sol} summarises the results of a simulation: for each sample size $n$, 800 independent samples of size $n$ were drawn from an exponential distribution with mean $3$, and the sample mean was recorded for each run.

\begin{table}[h]
\centering
\caption{Behaviour of Sample Means}\label{tab:b-sim-sol}
\begin{tabular}{@{}rcc@{}}
\toprule
\textbf{Sample size $n$} & \textbf{Average of the 800 sample means} & \textbf{SD of the 800 sample means}\\
\midrule
10  & 3.06 & 0.94 \\
40  & 3.01 & 0.47 \\
160 & 3.00 & 0.24 \\
\bottomrule
\end{tabular}
\end{table}

\begin{table}[h]
\centering
\caption{Selected QQ-Plot Quantiles}\label{tab:b-qq-sol}
\begin{tabular}{@{}rc@{}}
\toprule
\textbf{Normal quantile} & \textbf{Sample quantile}\\
\midrule
-2.0 & -1.4 \\
-1.0 & -0.8 \\
 0.0 &  0.1 \\
 1.0 &  0.9 \\
 2.0 &  1.6 \\
\bottomrule
\end{tabular}
\end{table}

\newpage

% =============================
% Solutions and grading notes
% =============================
\begin{center}
{\LARGE Solutions and Grading Notes (Version B)}\\[0.5em]
\end{center}

\section*{Question 1 — Probability Foundations (25 points)}

\begin{enumerate}
  \item \textbf{Probability space (6 pts)}\\
  Sample space: $\Omega = \{(i,j): i,j\in\{1,2,3,4,5,6\}\}$, consisting of 36 ordered pairs.\\
  Natural $\sigma$-algebra: $\mathcal{F} = 2^{\Omega}$.\\
  Probability measure: For the fair die, $P(\omega) = \tfrac{1}{36}$ for each $\omega\in\Omega$.\\
  Partial credit: correct $\Omega$ (2 pts), acceptable $\mathcal{F}$ (2 pts), correct probabilities summing to 1 (2 pts).

  \item \textbf{Event probabilities (6 pts)}\\
  $A=\{(4,6),(5,5),(5,6),(6,4),(6,5),(6,6)\}$, so $P(A)=\tfrac{6}{36}=\tfrac{1}{6}$.\\
  $B=\{(5,j),(6,j): j\in\{1,2,3,4,5,6\}\}$, so $P(B)=\tfrac{12}{36}=\tfrac{1}{3}$.\\
  $A\cap B=\{(5,5),(5,6),(6,4),(6,5),(6,6)\}$, so $P(A\cap B)=\tfrac{5}{36}$.\\
  Award 3 pts for correct event descriptions, 3 pts for the resulting probabilities.

  \item \textbf{Independence check (5 pts)}\\
  $P(A)P(B)=\tfrac{1}{6}\cdot\tfrac{1}{3}=\tfrac{1}{18}=\tfrac{2}{36}\neq\tfrac{5}{36}=P(A\cap B)$. Therefore $A$ and $B$ are \emph{not} independent.\\
  Credit: 3 pts for the product comparison, 2 pts for the explicit conclusion.

  \item \textbf{Biased die analysis (8 pts)}\\
  With outcomes 1--5 equally likely with total probability $1-q$, each has probability $\tfrac{1-q}{5}$; outcome 6 has probability $q$. Event $B$ has
  $P_q(B)=\tfrac{1-q}{5}+q=\tfrac{1+4q}{5}$. Setting $P_q(B)=\tfrac{1}{3}$ yields $q=\tfrac{1}{6}$, which lies on the boundary of the given range ($0<q<\tfrac{1}{6}$), so there is no valid solution strictly within the range.\\
  Award 5 pts for correct $P_q(B)$, 3 pts for solving and noting the constraint.
\end{enumerate}

\section*{Question 2 — Bayes' Rule in Fraud Detection (25 points)}

Let $F$ denote ``fraudulent'' and $+$ denote ``flagged''. Given $P(F)=\tfrac{1}{50}$, $P(+\mid F)=\tfrac{4}{5}$, and $P(-\mid F^c)=\tfrac{39}{40}$, so $P(+\mid F^c)=\tfrac{1}{40}$.

\begin{enumerate}
  \item \textbf{Probability of a flag (7 pts)}\\
  $P(+) = P(+\mid F)P(F) + P(+\mid F^c)P(F^c) = \tfrac{4}{5}\cdot\tfrac{1}{50} + \tfrac{1}{40}\cdot\tfrac{49}{50} = \tfrac{81}{2000}$.

  \item \textbf{Positive predictive value (8 pts)}\\
  $P(F\mid +) = \dfrac{\tfrac{4}{5}\cdot\tfrac{1}{50}}{\tfrac{81}{2000}} = \dfrac{\tfrac{16}{1000}}{\tfrac{81}{2000}} = \dfrac{32}{81}$.

  \item \textbf{Posterior after a negative (6 pts)}\\
  $P(-) = P(-\mid F^c)P(F^c) + P(-\mid F)P(F) = \tfrac{39}{40}\cdot\tfrac{49}{50} + \tfrac{1}{5}\cdot\tfrac{1}{50} = \tfrac{1919}{2000}$. Hence $P(F\mid -) = \dfrac{\tfrac{1}{5}\cdot\tfrac{1}{50}}{\tfrac{1919}{2000}} = \dfrac{8}{1919}$.

  \item \textbf{General formula and interpretation (4 pts)}\\
  With $\pi=P(F)$, $s=P(+\mid F)$ and $c=P(-\mid F^c)$: $\displaystyle P(F\mid +) = \frac{s\pi}{s\pi + (1-c)(1-\pi)}$. Increasing $c$ decreases $(1-c)$, reducing the denominator and thereby increasing $P(F\mid +)$.
\end{enumerate}

\section*{Question 3 — Uniform Waiting-Time Model (25 points)}

\begin{enumerate}
  \item \textbf{CDF (5 pts)}\\
  For $T\sim\mathrm{Uniform}(0,8)$, $\displaystyle F_T(t)=\begin{cases}0,& t<0,\\ t/8, & 0\le t\le 8,\\ 1, & t>8.\end{cases}$

  \item \textbf{Probability of an interval (5 pts)}\\
  $P(2\le T\le 6)=F_T(6)-F_T(2)=6/8-2/8=1/2$.

  \item \textbf{Expectation and variance via integration (9 pts)}\\
  Density: $f_T(t)=1/8$ for $t\in[0,8]$. $\displaystyle E[T]=\int_0^8 t\cdot\tfrac{1}{8}\,dt=4$. $\displaystyle E[T^2]=\int_0^8 t^2\cdot\tfrac{1}{8}\,dt=\tfrac{64}{3}$. So $\operatorname{Var}(T)=\tfrac{64}{3}-16=\tfrac{16}{3}$.

  \item \textbf{Affine transformation (6 pts)}\\
  $E[Y]=E[2T+5]=2E[T]+5=13$, \quad $\operatorname{Var}(Y)=2^2\operatorname{Var}(T)=\tfrac{64}{3}$.
\end{enumerate}

\section*{Question 4 — Interpreting Summaries and Convergence (25 points)}

\begin{enumerate}
  \item \textbf{Law of Large Numbers interpretation (9 pts)}\\
  The average of the sample means in Table~\ref{tab:b-sim-sol} stays close to the true mean 3: 3.06 for $n=10$, 3.01 for $n=40$, and 3.00 for $n=160$. As $n$ increases, the averages move closer to 3 and fluctuate less, illustrating that sample means converge toward the population mean, as predicted by the Law of Large Numbers.

  \item \textbf{Central Limit Theorem interpretation (8 pts)}\\
  The standard deviation of the 800 sample means decreases roughly by a factor close to $1/\sqrt{n}$: 0.94 for $n=10$, 0.47 for $n=40$, and 0.24 for $n=160$. This shrinking variability reflects the CLT prediction that $\bar X_n$ becomes more concentrated and approximately normal with variance $\sigma^2/n$ as $n$ grows. Full credit requires explicitly linking the observed reductions to the CLT scaling.

  \item \textbf{QQ-plot assessment and recommendation (8 pts)}\\
  Sample quantiles are less extreme in both tails than the normal quantiles (e.g., $-1.4>-2$ and $1.6<2$), indicating a QQ-plot that curves toward the center with lighter tails than normal. Therefore the dataset appears to have lighter tails than a standard normal distribution. Suggested actions: check for truncation, use a bounded distribution model, or assess whether normal-based inference remains appropriate given the lighter tails.
\end{enumerate}

\end{document}
